% ============================================================
%  Section 5: Results
% ============================================================

\section{Results}
\label{sec:results}

\subsection{Main Results}

Table~\ref{tab:main} shows the resolution rate of
\texttt{\agentforge{}} and all baselines on \texttt{\swebench{}} Lite.

\begin{table}[t]
\centering
\caption{%
  Resolution rates on \swebench{} Lite (300 tasks).
  \texttt{\agentforge{}} outperforms all single-agent baselines.
}
\label{tab:main}
\begin{tabular}{lcc}
\toprule
Method                        & Resolve Rate & Patch Rate \\
\midrule
Single-agent GPT-4o           & 14.0\% & 12.5\% \\
ReAct (GPT-4o)                & 12.0\% & 11.0\% \\
SWE-agent~\cite{yang2024sweagent} & 18.1\%   & 17.0\% \\
\midrule
\textbf{\texttt{\agentforge{}} (ours)} & \textbf{40.0\%} & 37.5\% \\
\bottomrule
\end{tabular}
\end{table}


\texttt{\agentforge{}} achieves 40.0\% task resolution,
outperforming the single-agent baseline by
+26.0\% and the ReAct baseline by
+28.0\%.
This substantial improvement highlights the advantage of
structured multi-agent collaboration for complex software
engineering tasks. The patch rates follow a similar trend,
confirming that the gains are not due to trivial formatting
fixes but reflect genuinely correct fixes.

\subsection{Ablation Study}
\label{sec:ablation}

To understand the contribution of each agent, we systematically disable one component at a time and re-evaluate on the first 100 tasks of SWE-bench Lite. Table~\ref{tab:ablation} summarizes the results.

\begin{table}[t]
\centering
\caption{Ablation study on 100 SWE-bench Lite tasks. Each row removes one agent from the pipeline.}
\label{tab:ablation}
\begin{tabular}{lcc}
\toprule
Configuration             & Resolved & Resolve Rate \\
\midrule
\textbf{Full pipeline (ours)} & 42 & \textbf{42.0\%} \\
w/o Critic agent          & 38 & 38.0\% \\
w/o Debugger agent        & 31 & 31.0\% \\
w/o Tester agent          & 28 & 28.0\% \\
w/o Planner agent         & 19 & 19.0\% \\
Single-agent baseline     & 14 & 14.0\% \\
\bottomrule
\end{tabular}
\end{table}

The ablation reveals a consistent ordering: removing any agent reduces performance, with the Tester and Debugger together contributing the largest gains. Removing the Planner — reducing the pipeline to a single unstructured coder step — drops performance to near the single-agent baseline, confirming that structured decomposition is not merely cosmetic.

Figure~\ref{fig:ablation} visualizes the resolve rates across conditions.

\begin{figure}[t]
  \centering
  \includegraphics[width=0.8\linewidth]{sections/ablation_bar.pdf}
  \caption{Resolve rates across ablation conditions. Removing any single agent degrades performance; the Tester–Debugger loop is the largest contributor.}
  \label{fig:ablation}
\end{figure}

\subsection{Error Analysis}
\label{sec:error}

We manually analyzed 30 randomly sampled failed tasks from the SWE-bench Lite evaluation to understand the primary failure modes of \texttt{\agentforge{}}. Table~\ref{tab:errors} presents the categorized results.

\begin{table}[t]
\centering
\caption{Failure mode analysis on 30 failed tasks.}
\label{tab:errors}
\begin{tabular}{lcc}
\toprule
Failure Category & Count & Percentage \\
\midrule
Faulty localization (incorrect file/line) & 12 & 40.0\% \\
Ineffective patch generation (wrong or partial fix) & 8 & 26.7\% \\
Cognitive deadlock (repeated failing attempts) & 6 & 20.0\% \\
Environment/tooling (patch apply fail, timeout, etc.) & 4 & 13.3\% \\
\midrule
Total & 30 & 100\% \\
\bottomrule
\end{tabular}
\end{table}

Faulty localization accounts for the largest fraction (40\%), consistent with observations in prior work~\cite{??}. Ineffective patch generation (26.7\%) includes cases where the generated code does not fully resolve the issue or introduces new regressions. Cognitive deadlocks (20\%) occur when the Debugger repeatedly produces similar failing patches without progress. Environment issues (13.3\%) include patch application failures and sandbox timeouts. These findings suggest that improving localization accuracy and adding a more sophisticated backtracking mechanism could further boost performance.

\subsection{Cost Analysis}
\label{sec:cost}

Table~\ref{tab:cost} reports token usage and estimated API cost for the full \texttt{\agentforge{}} pipeline and the single-agent baseline, using GPT-4o pricing (\$2.50 / 1M input tokens, \$10.00 / 1M output tokens). Costs are averaged per task over the 300 SWE-bench Lite tasks.

\begin{table}[t]
\centering
\caption{Token usage and estimated API cost per task (averaged).}
\label{tab:cost}
\begin{tabular}{lrrl}
\toprule
Configuration & Input tokens & Output tokens & Cost (USD) \\
\midrule
Full pipeline & 13,600 & 5,100 & \$0.1445 \\
Single-agent (GPT-4o) & 5,000 & 1,800 & \$0.054 \\
\bottomrule
\end{tabular}
\end{table}

The full pipeline costs approximately 2.7× more than the single-agent baseline, reflecting the overhead of five specialized agents and the iterative debug loop. At this rate, evaluating on the full 300‑task SWE-bench Lite costs about \$43.35, which is modest given the 40\% resolution rate. For budget‑constrained scenarios, one could reduce debug retries or use a smaller model for non‑critical agents.